% ---
% title: the RCGM framework
% date: 2025/11/12
% tag: Diffusion
% ---
\definecolor{mdc1}{HTML}{FFFFFF}
\definecolor{mdc2}{HTML}{0000DD}
\definecolor{mdc3}{HTML}{DD0000}
\section{Distillation}

蒸馏就是拿已有的扩散模型 \(\theta\)，减少采样步数，提高生成效率。常见做法有 DDIM、Reflow 等。

\section{Consistency Model}

这里的 noise schedule 稍有不同，\(x_t=x_0+tx_1\)。常规 one-step 蒸馏的训练目标是：

\[
\min \mathbf{E}\|f_{\theta}(x_1,1)-x_0\|
\]

一致性模型利用的是这样一个分解：

\[
\|f_{\theta^*}(x_1,1)-x_0\|=\|\large\sum_{k=1}^{n}[f_{\theta^*}(x_{t_k},t_k)-f_{\theta^*}(x_{t_{k-1}},t_{k-1})]\|
\le \large\sum_{k=1}^{n}\|[f_{\theta^*}(x_{t_k},t_k)-f_{\theta^*}(x_{t_{k-1}},t_{k-1})]\|
\]

边界条件是 \(t=0\) 时 \(f_\theta(x_0,0)=x_0\)，于是：

\[
f_\theta(x_{t_1},t_1) \to x_0 \\
f_\theta(x_{t_k},t_k) \to f_\theta(x_{t_{k-1}},t_{k-1})
\]

最终也得到：

\[
f_\theta(x_1,1) \approx x_0
\]

所以如果 Reflow 算一种蒸馏，Consistency Model 既可看成蒸馏，也可看成自训练。

\section{Short-cut Model}

\subsection{InstraFlow}

InstaFlow 和 Rectified Flow（Rect Flow）的關係可以概括成一句話：Rectified Flow 是技術原理，InstaFlow 是用這個技術做出來的具體模型。

可以把它們想成「渦輪增壓技術（Turbo）」和「某款特定的渦輪跑車」的關係。

\subsubsection{Rectified Flow (Rect Flow)：核心技術}

這是一個通用的數學框架和訓練方法（由 Liu et al. 提出）。

它的目標是把擴散模型原本彎彎曲曲的生成路徑（trajectory）拉直。原本的路徑很彎，採樣時得走很多小碎步（例如 20 到 50 步）才不會走歪。Rectified Flow 用一種叫 Reflow（重流）的反覆訓練，讓雜訊到圖片的路徑變成直線。變直之後通常還是個少步數模型，例如從 50 步減到 4 到 8 步，並不是專門為一步生成設計的。

\subsubsection{InstaFlow：具體模型}

這是一篇論文（基於 Stable Diffusion 改進）提出的一個具體模型名稱。

它的目標是一步生成（one-step generation）。做法是在 Rectified Flow 之上再加一次蒸餾：先用 Reflow 把 Stable Diffusion 的路徑拉直，這時可能還要 2 到 3 步；再蒸餾一次，強迫模型在 1 步內走完整條直線。結果是一個專門的 1-step 模型，約 0.1 秒一張圖。

\hrule

\subsubsection{對比表}

\begin{center}
\begin{tabular}{lll}
\toprule
\textbf{特性} & \textbf{Rectified Flow (Rect Flow)} & \textbf{InstaFlow} \\
\midrule
\textbf{定義} & 一種\textbf{訓練方法 / 數學框架} & 一個\textbf{具體的模型 / 論文成果} \\
\textbf{數學本質} & 讓生成軌跡變直 (Straightening) & Rectified Flow + 蒸餾 (Distillation) \\
\textbf{主要目的} & 減少採樣步數，提高 ODE 穩定性 & \textbf{實現 1 步生成 (One-Step)} \\
\textbf{生成步數} & 通常 2\textasciitilde{}8 步 (視 Reflow 次數而定) & \textbf{嚴格 1 步} \\
\textbf{關係} & 它是 InstaFlow 的地基 & 它是 Rectified Flow 的具體應用 \\
\bottomrule
\end{tabular}
\end{center}

所以在論文或技術文檔裡看到 Rect Flow，通常是在講「把線拉直」這個數學技巧；看到 InstaFlow，則是指那個用 Rect Flow 做出來、能一步出圖的 Stable Diffusion 改版模型。

\subsubsection{InstaFlow Step 2 - Distillation}

InstaFlow 第二步的蒸餾能一步出圖，關鍵不在蒸餾算法本身，而在於 Step 1 的 Reflow 已經把問題簡化了。這是「前人栽樹（拉直），後人乘涼（蒸餾）」的過程。

\paragraph{核心邏輯：為什麼能蒸餾？}

傳統擴散模型（如 SD 1.5/2.0）裡，從雜訊 \(z\) 到圖像 \(x\) 的映射是一條很彎的非線性路徑。要讓模型一步預測終點，難度很高（underfitting，學不起來）。經過 Step 1 的 Reflow 之後，傳輸軌跡變成了直線。

軌跡是直線時，\(x\) 和 \(z\) 的關係退化成簡單的線性關係（或接近線性），而\colorbox{mdc2}{\textcolor{mdc1}{神經網絡非常擅長擬合線性或平滑的映射}}。所以 InstaFlow 不用像 Consistency Distillation (CD) 那樣加一堆自洽性約束，直接用最樸素的直接蒸餾（direct distillation）。

\paragraph{具體實現步驟}

InstaFlow 的蒸餾是一個 teacher-student 架構。Teacher 模型（\(T\)）是 Step 1 訓練好的 Reflowed Model，通常叫 2-Rectified Flow，生成的軌跡已經很直。Student 模型（\(S\)）就是要訓練的 InstaFlow（one-step model），初始權重通常複製自 Teacher。

\paragraph{訓練流程 (Algorithm)}

每次訓練迭代：

\begin{enumerate}
  \item 採樣 (Sampling)：
  
  隨機抽取一個高斯雜訊 \(z \sim \mathcal{N}(0, I)\) 和對應的文本提示 \(c\)。
  \item 老師生成\colorbox{mdc3}{\textcolor{mdc1}{"偽真值" (Teacher Generation)}}：
  
  使用 Teacher 模型，通過常微分方程求解器 (ODE Solver，例如 Euler method)，用較多的步數（例如 25-50 步）從 \(z\) 生成一張圖像 \(x_{teacher}\)。
  
  \[
  x_{teacher} = \text{ODESolve}(T, z, t=1 \rightarrow 0, c)
  \]
  
  注意：這裡的 \(x_{teacher}\) 被視為 Ground Truth。
  \item 學生預測 (Student Prediction)：
  
  將同樣的雜訊 \(z\) 和提示 \(c\) 輸入給 Student 模型，要求它 一步 (One-step) 直接輸出圖像 \(x_{student}\)。
  
  \[
  x_{student} = S(z, c)
  \]
  \item 計算損失 (Loss Calculation)：
  
  直接計算兩者之間的 L2 損失 (或 LPIPS 感知損失)：
  
  \[
  \mathcal{L} = \| x_{student} - x_{teacher} \|^2
  \]
  \item 反向傳播：更新 Student 模型的參數。
\end{enumerate}

\paragraph{與傳統蒸餾 (Progressive Distillation) 的區別}

Salimans 的 Progressive Distillation 走步數減半（64 步 \(\to\) 32 步 \(\to\) 16 步……）：原始軌跡是彎的，不能直接跳到底，必須慢慢合併步數。InstaFlow 是直接折疊（direct collapse）：軌跡已經是直線，理論上

\[
\text{位移} = \text{平均速度} \times \text{時間}
\]

直線運動的瞬時速度等於平均速度，所以 Student 只要學會預測這個恆定速度（或直接預測終點），就能一步到位。

\paragraph{優化目標}

InstaFlow 的蒸餾可以寫成這樣一個優化問題：

\[
\min_\theta \mathbb{E}_{z \sim \mathcal{N}(0, I)} \left[ \| S_\theta(z) - \text{Reflowedode}(z) \|^2 \right]
\]

其中 \(\text{Reflowedode}(z)\) 是那個已經被「拉直」過的老師模型的輸出。

蒸餾本身沒什麼花俏的地方，起作用的是老師模型（\(T\)）已經把路鋪直了；學生模型（\(S\)）只要沿直線走，不用學轉彎。

\subsection{Distribution Matching Distillation}

DMD（Distribution Matching Distillation）要解決的是 SDS（Score Distillation Sampling）的缺陷，把它變成一個能穩定訓練的單步生成框架。

\hrule

\subsubsection{研究背景 (Research Background)}

要解決的問題：把多步的 diffusion model（老師）蒸餾成單步的 generator（學生），畫質不能掉。

\colorbox{mdc3}{\textcolor{mdc1}{蒸餾的兩種路線：基於採樣軌跡的，和基於採樣分佈的}}

傳統蒸餾（如 Progressive Distillation）大多依賴 MSE loss（\(L_2\) 範數）。給定一個輸入噪聲 \(z\)，老師模型對應的真實分佈 \(p(x|z)\) 往往是多模態的：一個模糊的輪廓，既可能是貓，也可能是狗。MSE 的最優解是這些模態的期望值，而期望值通常落在高概率密度區域的中間（低概率區域），表現在圖像上就是所有可能結果的疊加，也就是模糊。

DreamFusion 的 SDS loss 改用 score function 引導生成，避開了 MSE 的模糊問題，但它忽略了生成分佈自身的熵項（fake distribution 的 score），生成的圖像飽和度過高、紋理怪異（Janus problem）。

\subsubsection{研究動機 (Research Motivation)}

DMD 的假設是：學生模型 \(G_\theta(z)\) 的輸出不必和老師的軌跡逐像素對齊（pixel-wise alignment），只要學生生成的圖像集合 \(p_{fake}\) 和真實數據分佈 \(p_{real}\) 在統計上分不出來就行。

數學上就是最小化兩個分佈的 KL 散度（Kullback-Leibler divergence）：

\[
\min_\theta D_{KL}(p_{fake} \| p_{real}) = \mathbb{E}_{x \sim p_{fake}} \left[ \log p_{fake}(x) - \log p_{real}(x) \right]
\]

難點在於梯度算不出來：\(p_{fake}(x)\) 是學生網絡 \(G_\theta\) 定義的隱式分佈，寫不出解析式，\(p_{real}(x)\) 也拿不到。DMD 要做的就是找一個可計算、穩定的梯度估計來優化這個 KL 散度。

\subsubsection{研究方法 (Research Method)}

總損失分成兩部分：\(\mathcal{L}_{total} = \mathcal{L}_{dist} + \lambda \mathcal{L}_{reg}\)。重點看這兩部分怎麼推出來。

\paragraph{第一部分：分佈匹配損失 (Distribution Matching Loss)}

這部分用預訓練的擴散模型（老師）來估計 KL 散度的梯度。

A. 梯度的推導

我們想優化 generator 參數 \(\theta\)。對於 KL 散度 \(J = D_{KL}(p_\theta \| p_{real})\)，其梯度可以推導為（省略部分鏈式法則細節）：

\[
\nabla_\theta J \approx \mathbb{E}_{z \sim \mathcal{N}(0,I)} \left[ \nabla_\theta G(z) \cdot (\underbrace{\nabla_x \log p_{fake}(G(z))}_{\text{Fake Score}} - \underbrace{\nabla_x \log p_{real}(G(z))}_{\text{Real Score}}) \right]
\]

這是一個力學平衡：real score 把生成點 \(x\) 拉向真實數據的高密度區域，fake score（熵項）則把 \(x\) 推離當前生成的擁擠區域，防止 mode collapse。

B. 怎麼算 Real Score？

這是擴散模型的強項。根據 Tweedie's Formula 或 Score Matching 的定義，預訓練的老師模型 \(\epsilon_\phi(x_t, t)\) 本質上就是真實分佈的 Score 估計器（在噪聲等級 \(t\) 下）：

\[
\nabla_{x_t} \log p_{real}(x_t) \propto - \epsilon_\phi(x_t, t)
\]

所以隨機採樣一個時間步 \(t\)，把學生生成的圖 \(x_{fake}\) 加噪得到 \(x_t\)，丟給老師算 \(\epsilon\)，就得到真實分數的方向。

C. 怎麼算 Fake Score？（DMD 的關鍵 trick）

\(p_{fake}\) 是學生自己定義的，score 很難算。DMD 的近似是把 fake score 當成高斯分佈的 score：假設在加噪後的 \(x_t\) 空間裡，學生的分佈 \(p_{fake}(x_t)\) 可以局部近似成中心在 \(\mu_{fake}\) 的高斯，於是 fake score 變成：

\[
\nabla_{x_t} \log p_{fake}(x_t) \approx \frac{x_t - \mu_{fake}}{\sigma_t^2}
\]

而 \(\frac{x_t - \mu_{fake}}{\sigma_t}\) 就是加進去的噪聲 \(\epsilon\)。

D. 最終的梯度形式

把 B 和 C 合起來，DMD 實際在算：

\[
\nabla_\theta \mathcal{L}_{dist} \propto \mathbb{E}_{t, \epsilon} \left[ \nabla_\theta G(z) \cdot (\epsilon_\phi(G(z)_t, t) - \epsilon) \right]
\]

\(\epsilon_\phi(\dots)\) 是老師認為這個噪聲該長什麼樣，代表真實圖像結構；\(\epsilon\) 是實際加進去的噪聲，代表學生當前的分佈結構。

如果老師預測的噪聲和實際噪聲對不上，說明這張圖不夠真，梯度就會推著 \(G(z)\) 去修正差異。這其實是 SDS loss 的修正版，也可以看成把擴散模型當成 GAN 的 discriminator。

\paragraph{第二部分：回歸損失 (Regression Loss)}

\[
\mathcal{L}_{reg} = \| G_\theta(z) - \mu_\phi(z, T_{step}) \|^2
\]

純粹優化 KL 很不穩定，容易產生幻覺或跟 prompt 對不上，所以作者留了一個弱約束：利用老師模型做一步或多步確定性採樣（例如用 DDIM 跑幾步），作為 \(y_{target}\)。這是一個 anchor（錨點）：細節可以自由發揮，由 dist loss 負責，但大體構圖和顏色（低頻信息）要跟著老師走。

\paragraph{訓練流程}

\begin{enumerate}
  \item 採樣噪聲 \(z\)。
  \item Student forward：\(x_{fake} = G_\theta(z)\)。
  \item Regression branch：計算 \(x_{fake}\) 和老師生成的粗糙軌跡之間的 MSE。
  \item Distribution branch：
  
  \begin{itemize}
    \item 隨機選一個 \(t\)。
    \item 給 \(x_{fake}\) 加噪得到 \(x_t\)。
    \item 讓老師模型看 \(x_t\)，吐出 \(\epsilon_{teacher}\)。
    \item 計算梯度差值 \((\epsilon_{teacher} - \epsilon_{noise})\) 並反向傳播。
  \end{itemize}
  \item 更新 \(\theta\)。
\end{enumerate}

\subsubsection{幾點結論}

FP16 訓練的可行性是一個工程上的發現。之前的蒸餾方法（如 CD）對數值精度很敏感，通常要 FP32，DMD 證明合理的 loss 設計可以讓半精度訓練穩定下來，這對資源有限的實驗室很有用。

DMD 還把「內容」和「質感」拆成了兩個正交任務：MSE 負責低頻內容（content/structure），score distillation 負責高頻紋理（texture/realism）。這個拆法對後面設計新的 loss function 有啟發。

GAN 與 diffusion 還有一個統一視角。DMD 相當於把 diffusion model 變成一個「完美的 discriminator」：GAN 的判別器容易訓練崩潰，預訓練 diffusion model 給出的梯度場則平滑、穩定。照這個思路，生成模型可能會走向 generator（one-step）加 diffusion-based critic 的架構。

如果要復現或改進這個工作，重點看 fake score 的近似。DMD 用高斯近似是一個簡化，如果能引入更精確的 entropy 估計，例如加一個額外的 head 預測 entropy，大概會是下一篇 paper 的點。
